\paragraph*{04.05.01.01. Erfolgskriterien anerkennen, priorisieren und überprüfen}
\kcilabel{para:04.05.01.01_ErfolgskriterienAnerkennenPriorisierenUndÜberprüfen}{04.05.01.01}
\label{para:04.05.01.01_ErfolgskriterienAnerkennenPriorisierenUndÜberprüfen}

% Task
\par In dem langlaufenden, komplexen Projektumfeld mit hohen Abhängigkeiten, zahlreichen Stakeholdern und parallel laufendem Betrieb von \gls{r3} und \gls{s4} war es meine Aufgabe, \gls{e2e}-übergreifende Erfolgskriterien aktiv zu steuern. Ich konzentrierte mich nicht nur auf die technische Steuerung der \gls{ricefw}-Liefergegenstände, sondern lenkte den Fokus auch auf die fachlichen, technischen und organisatorischen Ziele des Programs.

% Action
\par In enger Abstimmung mit den Stakeholdern definierte, priorisierte und überprüfte ich die Erfolgskriterien entlang der \gls{e2e}-Prozesskette. Als \gls{pm} integrierte ich diese Kriterien in die tägliche Projektarbeit und stellte sicher, dass sie von allen Beteiligten verstanden und verfolgt wurden. Dazu gehörte auch der bewusste Verzicht auf den \gls{aif}-Ansatz zugunsten von \gls{cc} als gezielte \gls{td}, um die priorisierten Kriterien zu erfüllen. Die wichtigsten Kriterien waren: vollständige und fehlerfreie Abdeckung der fachlichen Kernprozesse der Buchungskreise, Einhaltung technischer Qualitätsstandards für \gls{ricefw}-Objekte, Minimierung von Störungen in der \gls{hypc}-Phase sowie organisatorische Absicherung des Parallelbetriebs.

% Result
\par Die konsequente Fokussierung auf priorisierte Erfolgskriterien sicherte den technischen Erfolg von \gls{cod}, \gls{hypc} und der Teilübergabe des \gls{s4} in den Linienbetrieb sowie die Erfüllung fachlicher Stakeholder-Anforderungen. Der laufende Betrieb blieb während der Projektlaufzeit stabil. Die kontinuierliche Überprüfung und Anpassung der Kriterien ermöglichte flexible Reaktionen auf Veränderungen und hielt das Projekt auf die wichtigsten Ziele ausgerichtet.

% Reflect
\par Die Erfahrung zeigt, dass klare Definition und Priorisierung von Erfolgskriterien in komplexen Projekten entscheidend sind. Erfolgskriterien müssen integraler Bestandteil von Planung und Durchführung sein, damit alle Aktivitäten auf die übergeordneten Ziele ausgerichtet sind. Die fortlaufende Überprüfung und Anpassung der Kriterien ist erforderlich, um auf Veränderungen im Projektumfeld zu reagieren und den Fokus zu halten.